\documentclass[11pt]{article}
\usepackage[margin=0.9in]{geometry}
\usepackage[T1]{fontenc}
\usepackage{times}
\usepackage{microtype}
\usepackage[hidelinks]{hyperref}
\usepackage{parskip}
\pagestyle{empty}
\setlength{\parskip}{7pt}

\begin{document}

\enlargethispage{1.5\baselineskip}

\begin{flushleft}
\textbf{Yichao Jin, Ph.D.}\\
School of Economic, Political and Policy Sciences\\
The University of Texas at Dallas\\
\href{mailto:Yichao.Jin@UTDallas.edu}{Yichao.Jin@UTDallas.edu} \,|\, \href{https://yichao2022.github.io}{yichao2022.github.io}
\end{flushleft}

\vspace{14pt}

August 17, 2026

\vspace{14pt}

Professor Jagpreet (Jag) Chhatwal, PhD\\
Director, Center for Health Technology Assessment\\
Mass General Brigham\\
Harvard Medical School

\vspace{14pt}

\textbf{Re:} Postdoctoral Fellowship Application

\vspace{14pt}

Dear Professor Chhatwal,

I am writing to apply for the Postdoctoral Fellow position at the Center for Health Technology Assessment, Mass General Brigham and Harvard Medical School. I recently completed my Ph.D. in Public Policy and Political Economy at the University of Texas at Dallas (May 2026), where I built a research program at the intersection of health preference measurement, decision-analytic modeling, and generative AI for health outcomes research.

My methodological home is discrete choice experiments combined with structural econometrics. My dissertation, based on a discrete choice experiment of over 600 respondents in Wuhan, quantified how waiting time and institutional trust jointly shape vaccination uptake; the resulting paper, ``Waiting Time as a Behavioral Barrier to Vaccination Uptake: Nonlinear Heterogeneity by Institutional Trust,'' is under review at \textit{Economic Analysis and Policy}. I have also developed a \textbf{Behavioral Digital Twin (BDT)} framework that integrates DCE data with causally-constrained LLM agents to generate synthetic policy simulations. My working paper on this method, ``Empirical-Frontier Regularization for LLM Synthetic Agents: A Preference-Anchored Framework'' (targeted at \textit{Value in Health}), demonstrates an anchor-based calibration that reduces waiting-time monotonicity violations from \textbf{62.5\% to 4.2\%} and raises Spearman rank correlation from \textbf{$-$0.024 to 0.952}. A companion paper, ``Value Contamination in Policy Simulation,'' audits how LLM policy-value orientations bias predicted policy effects. I see these tools---preference measurement, simulation, and the emerging role of generative AI in HEOR---as directly aligned with CHTA's agenda in decision science and health technology assessment.

CHTA's work has been a reference point throughout my training. Your hepatitis C elimination models translated disease modeling into federal policy; your COVID-19 modeling consortium provided real-time guidance to states and the CDC; and your recent HEALing Communities simulation shows how modeling can guide opioid policy at the level of individual communities. I would welcome the opportunity to contribute to that tradition, developing models that move from academic insight to clinical and policy decisions.

Beyond these projects, I bring advanced skills in Python, R, Stata, and MATLAB; experience with mixed logit, latent class, and SEIR modeling; and a strong interest in pursuing independent funding (e.g., an NIH K Award). I am flexible on start date and would be glad to begin before or at Fall 2026.

My CV is attached, and full drafts of both working papers are available on SSRN. Thank you for your time and consideration. I would be delighted to discuss how I might contribute to CHTA's research portfolio.

\vspace{12pt}

Sincerely,

\vspace{1pt}

Yichao Jin, Ph.D.

\end{document}